\scanpage{58}
scindage de $E_U-E_K$ se remonte en un scindage de $E_L-E_K$ (ou aussi, \deleted{par} en un scindage de $E_V-E_K$, si $V$ est un modèle $\subset U$)

On va mettre les conjectures précédentes (sous forme modélique, disons) sous une forme plus géométrique, en introduisant, en $\unsure{=|\overline m}$ temps qu'un revêtement universel \deleted{\unsure{de}} $\widetilde U$ de $U$, le normalisé \deleted{de} $X'$ de $X$ dans $\widetilde U$ (où $X=\widehat U$). (NB je m'abstiens de le noter $\widetilde X$, car que justement il n'est pas étale sur $X$\ldots). Notons que pour $i\in\bar I=\bar X-\bar U$, les \inserted{l'ens des} $L_i$ de $\bar\pi=\pi(\bar U)$ sont en corr. $1-1$ avec l'\deleted{ens} fibre $X'_i$ de $X'$ au dessus de $i$

\diagram{\begin{tikzcd}[ampersand replacement=\&,column sep=large,row sep=large]
X\arrow[r,no head]\arrow[d,phantom,"\cup" description]\&\bar X\arrow[r,no head]\arrow[d,phantom,"\cup" description]\&X'\arrow[d,phantom,"\cup" description]\\
U\arrow[r,no head]\arrow[d,no head]\&\bar U\arrow[r,no head]\arrow[d,no head]\&\widetilde U\\
K\arrow[r,no head]\&\bar K\&
\end{tikzcd}}

Donc $X'$ peut être considéré comme la réunion de $\widetilde U$, et de $X'|\bar I=$ ensemble des sous-groupes lacets de $\bar\pi$, qui apparaissent ainsi comme des ``pts à l'infini'' du rev. universel $\widetilde U$. D'ailleurs $E_U$ s'inter-
\margin{57}
